%=====================================================================
\chapter{The Mathematical Toolkit}\label{ch:math}
%=====================================================================
\locator{1}{Part I --- Foundations of Data}

\begin{outcomes}
By the end of this chapter you will be able to:
\begin{tight}
  \item \textbf{Represent} an observation as a vector and a dataset as a matrix, and state what each dimension means.
  \item \textbf{Compute} a dot product and a Euclidean distance by hand, and explain what each measures about two observations.
  \item \textbf{Interpret} a derivative as a slope and a gradient as a direction of steepest increase.
  \item \textbf{Trace} one step of gradient descent numerically.
  \item \textbf{Explain} why every model in this book is ultimately an optimisation problem.
\end{tight}
\end{outcomes}

\begin{prereq}
\chref{data} --- observations, features and measurement scales. Nothing beyond
school algebra is assumed. If the prerequisite self-check troubled you, work this
chapter slowly; everything later depends on it.
\end{prereq}

\begin{keyterms}
scalar $\bullet$ vector $\bullet$ matrix $\bullet$ dot product $\bullet$ norm
$\bullet$ Euclidean distance $\bullet$ cosine similarity $\bullet$ function
$\bullet$ derivative $\bullet$ slope $\bullet$ gradient $\bullet$ minimum
$\bullet$ learning rate $\bullet$ gradient descent
\end{keyterms}

\section{Why Mathematics, and How Much}

Every model in this book does the same two things: it \emph{represents} data as
numbers arranged in a particular shape, and it \emph{searches} for the arrangement
of its own parameters that makes its errors smallest. The first is linear algebra;
the second is calculus. That is the whole of the mathematics you need, and this
chapter covers it at the depth this course requires --- intuition, notation, and
the ability to do a small example by hand.

\begin{conceptbox}[The One-Sentence Summary of Machine Learning]
\emph{Put the data in a matrix, define a number that says how wrong you are, and
roll downhill until that number stops getting smaller.} Everything in Parts III
and IV is a variation on this sentence.
\end{conceptbox}

\section{Vectors: One Observation}

\begin{definitionbox}[Scalar, Vector, Matrix]
A \term{scalar} is a single number ($5$, $-2.3$). A \term{vector} is an ordered
list of numbers, written $x = [x_1, x_2, \dots, x_p]$ --- one observation with $p$
features. A \term{matrix} is a rectangular grid, written $X$ with $n$ rows and $p$
columns --- a whole dataset.
\end{definitionbox}

A student described by (attendance 92\%, quizzes 8, forum posts 12) is the vector
$x = [92, 8, 12]$. That is not an analogy; it is literally how the data is stored
and how every algorithm sees it.

\dsfig{%
\begin{tikzpicture}[font=\scriptsize,
  cel/.style={draw=gray!50, fill=white, minimum width=1.25cm, minimum height=0.5cm},
  hcel/.style={draw=primary!60, fill=primary!12, text=primary!50!black,
               font=\scriptsize\bfseries, minimum width=1.25cm, minimum height=0.5cm},
  rcel/.style={draw=accent!60, fill=accent!10, minimum width=1.25cm, minimum height=0.5cm}]

% column headers
\foreach \c/\t in {0/attend, 1/quizzes, 2/posts}
  \node[hcel] at (\c*1.25,0) {\t};
% rows
\foreach \r/\a/\b/\c in {0/92/8/12, 1/45/3/0, 2/78/6/4}{
  \pgfmathsetmacro{\yy}{-0.5-\r*0.5}
  \ifnum\r=0
    \node[rcel] at (0,\yy) {\a}; \node[rcel] at (1.25,\yy) {\b}; \node[rcel] at (2.5,\yy) {\c};
  \else
    \node[cel] at (0,\yy) {\a}; \node[cel] at (1.25,\yy) {\b}; \node[cel] at (2.5,\yy) {\c};
  \fi}

\node[font=\scriptsize, text=accent, anchor=west] at (3.0,-0.5)
     {$\leftarrow$ one \textbf{observation} = one vector $x_1=[92,8,12]$};
\node[font=\scriptsize, text=primary!60!black, anchor=west] at (3.0,0.0)
     {$\uparrow$ one \textbf{feature} = one column};
\node[font=\scriptsize, text=gray!55!black, anchor=west] at (3.0,-1.35)
     {the whole grid = the \textbf{matrix} $X$, here $n=3$, $p=3$};

% dimension braces
\draw[decorate, decoration={brace, amplitude=4pt, mirror}, gray!60]
      (-0.72,-1.78) -- (3.12,-1.78) node[midway, below=4pt, font=\scriptsize] {$p=3$ features};
\draw[decorate, decoration={brace, amplitude=4pt}, gray!60]
      (-0.78,-1.75) -- (-0.78,-0.25) node[midway, left=4pt, font=\scriptsize] {$n=3$};
\end{tikzpicture}}%
{A dataset as a matrix. Every algorithm in this book consumes exactly this object: $n$
rows of $p$ numbers. Making real-world data fit this shape is the work of
\chref{wrangling}.}{matrix-anatomy}

\subsection{Length, Distance and Direction}

Three operations on vectors carry almost all the geometric intuition you need.

\begin{definitionbox}[Norm, Distance, Dot Product]
The \term{norm} (length) of $x$ is $\|x\| = \sqrt{x_1^2 + \dots + x_p^2}$.\\[2pt]
The \term{Euclidean distance} between $a$ and $b$ is
$d(a,b) = \sqrt{\sum_{j=1}^{p}(a_j - b_j)^2}$ --- how far apart two observations are.\\[2pt]
The \term{dot product} is $a \cdot b = \sum_{j=1}^{p} a_j b_j$ --- a single number
measuring how much two vectors point the same way.
\end{definitionbox}

\begin{worked}[Distance Between Two Students]
$a = [92, 8, 12]$ (engaged), $b = [45, 3, 0]$ (disengaged).

\smallskip
\textbf{Differences:} $92-45 = 47$, \quad $8-3 = 5$, \quad $12-0 = 12$.

\textbf{Squares:} $47^2 = 2209$, \quad $5^2 = 25$, \quad $12^2 = 144$.

\textbf{Sum:} $2209 + 25 + 144 = 2378$.

\textbf{Distance:} $d(a,b) = \sqrt{2378} \approx 48.8$.

\smallskip
\textbf{The lesson hiding in this number.} The attendance difference (47) supplied
$2209/2378 = 93\%$ of the total. Attendance dominates purely because it is measured
on a 0--100 scale while the others run 0--20. The distance is not measuring
similarity of students; it is measuring similarity of \emph{attendance}. This is
exactly why \chref{wrangling} insists on \term{scaling} features before any
distance-based method (k-NN in \chref{classification}, k-means in
\chref{unsupervised}) is used.
\end{worked}

\begin{conceptbox}[Cosine Similarity: Direction Without Magnitude]
Sometimes you care that two things are \emph{about the same}, not that they are the
same \emph{size}. Cosine similarity,
$\cos\theta = \dfrac{a \cdot b}{\|a\|\,\|b\|}$, ranges from $-1$ to $1$ and ignores
length entirely. It is the standard similarity measure for text
(\chref{nlp}), where a long document and a short one on the same topic should count
as similar.
\end{conceptbox}

\section{Functions, Slopes and Gradients}

\begin{definitionbox}[Derivative]
The \term{derivative} of $f$ at a point, written $f'(x)$ or $\frac{df}{dx}$, is the
\emph{slope} of $f$ there: how much $f$ changes per unit change in $x$. Where the
derivative is zero, the function is momentarily flat --- at a peak, a valley or a
plateau.
\end{definitionbox}

For this course you need three derivatives and one rule:

\begin{center}
$\dfrac{d}{dx}(c) = 0$, \qquad
$\dfrac{d}{dx}(x^2) = 2x$, \qquad
$\dfrac{d}{dx}(x^n) = n x^{n-1}$, \qquad
and derivatives add.
\end{center}

\dsfig{%
\begin{tikzpicture}
\begin{axis}[
  width=11.5cm, height=6.0cm,
  xlabel={\scriptsize parameter $w$}, ylabel={\scriptsize loss $L(w)$},
  xmin=-3.4, xmax=3.4, ymin=-0.6, ymax=9.5,
  axis lines=left, tick label style={font=\scriptsize},
  label style={font=\scriptsize}, legend style={font=\scriptsize, draw=gray!40,
  at={(0.99,0.99)}, anchor=north east},
  grid=major, grid style={gray!18},
]
\addplot[primary, line width=1.3pt, domain=-3.2:3.2, samples=120] {x^2};
\addlegendentry{$L(w)=w^2$}

% tangent at w = 2 (slope 4)
\addplot[accent, line width=1.1pt, dashed, domain=0.9:3.1] {4*x - 4};
\addlegendentry{slope at $w=2$ is $+4$}
\addplot[accent, only marks, mark=*, mark size=1.9pt] coordinates {(2,4)};

% tangent at w = -1.5 (slope -3)
\addplot[greenhl, line width=1.1pt, dashed, domain=-2.9:-0.2] {-3*x - 2.25};
\addlegendentry{slope at $w=-1.5$ is $-3$}
\addplot[greenhl, only marks, mark=*, mark size=1.9pt] coordinates {(-1.5,2.25)};

\addplot[purplehl, only marks, mark=*, mark size=2.4pt] coordinates {(0,0)};
\node[font=\scriptsize, purplehl, anchor=north] at (axis cs:0,-0.05) {minimum: slope $=0$};

\draw[-{Stealth[length=2mm]}, accent, line width=1pt]
     (axis cs:2,4) -- (axis cs:1.2,4);
\draw[-{Stealth[length=2mm]}, greenhl, line width=1pt]
     (axis cs:-1.5,2.25) -- (axis cs:-0.7,2.25);
\node[font=\scriptsize\itshape, text=gray!55!black, align=center] at (axis cs:0,7.2)
     {move \emph{against} the slope\\to go downhill};
\end{axis}
\end{tikzpicture}}%
{Slope as the guide to the minimum. Where the slope is positive the minimum lies to the
left; where it is negative, to the right. In both cases you move in the direction
\emph{opposite} the sign of the derivative --- which is the entire idea of gradient
descent.}{slope-minimum}

\begin{definitionbox}[Gradient]
When a function has several inputs, the \term{gradient} $\nabla L$ is the vector of
its partial derivatives --- one slope per parameter. It points in the direction of
\emph{steepest increase}, so $-\nabla L$ points downhill.
\end{definitionbox}

\section{Gradient Descent}

This is the algorithm that trains almost every model in Parts III and IV. It is
three lines long.

\begin{definitionbox}[Gradient Descent]
Repeat until the loss stops improving:
\[
w \;\leftarrow\; w \;-\; \alpha \, \nabla L(w)
\]
where $w$ are the parameters, $L$ is the loss, $\nabla L$ is its gradient, and
$\alpha$ is the \term{learning rate} --- how big a step to take.
\end{definitionbox}

\begin{worked}[Three Steps of Gradient Descent by Hand]
Minimise $L(w) = w^2$, starting at $w_0 = 3$, with learning rate $\alpha = 0.3$.
The gradient is $\nabla L = 2w$.

\smallskip
\textbf{Step 1.} $\nabla L(3) = 6$. \quad
$w_1 = 3 - 0.3 \times 6 = 3 - 1.8 = 1.2$. \quad $L = 1.44$.

\textbf{Step 2.} $\nabla L(1.2) = 2.4$. \quad
$w_2 = 1.2 - 0.3 \times 2.4 = 1.2 - 0.72 = 0.48$. \quad $L = 0.23$.

\textbf{Step 3.} $\nabla L(0.48) = 0.96$. \quad
$w_3 = 0.48 - 0.3 \times 0.96 = 0.48 - 0.288 = 0.192$. \quad $L = 0.037$.

\smallskip
The loss fell from $9$ to $0.037$ in three steps, and $w$ is converging on the true
minimum at $0$. Notice the steps get smaller automatically: near the minimum the
gradient is small, so the algorithm slows down as it arrives. Nobody had to program
that.

\smallskip
\textbf{Now try $\alpha = 1.2$.} $w_1 = 3 - 1.2 \times 6 = -4.2$; $w_2 = -4.2 -
1.2\times(-8.4) = 5.88$. The loss is \emph{rising}. Too large a learning rate
overshoots and diverges --- which is the subject of the next figure.
\end{worked}

\dsfig{%
\begin{tikzpicture}
\begin{groupplot}[
  group style={group size=3 by 1, horizontal sep=1.05cm},
  width=5.35cm, height=4.3cm,
  xmin=-4.6, xmax=4.6, ymin=-0.8, ymax=13,
  axis lines=left, tick label style={font=\tiny}, label style={font=\tiny},
  title style={font=\scriptsize\bfseries}, xlabel={$w$},
]
% --- too small
\nextgroupplot[title={\color{secondary}$\alpha$ too small}, ylabel={$L(w)$}]
\addplot[primary, line width=1pt, domain=-4.4:4.4, samples=80] {x^2};
\addplot[secondary, only marks, mark=*, mark size=1.5pt]
  coordinates {(3,9) (2.7,7.29) (2.43,5.9) (2.19,4.79) (1.97,3.87)};
\draw[-{Stealth[length=1.5mm]}, secondary] (axis cs:3,9) -- (axis cs:2.7,7.29);
\draw[-{Stealth[length=1.5mm]}, secondary] (axis cs:2.7,7.29) -- (axis cs:2.43,5.9);
\node[font=\tiny, text=secondary!70!black, align=center] at (axis cs:-1.6,10.5) {crawls;\\wastes time};

% --- just right
\nextgroupplot[title={\color{greenhl}$\alpha$ about right}]
\addplot[primary, line width=1pt, domain=-4.4:4.4, samples=80] {x^2};
\addplot[greenhl, only marks, mark=*, mark size=1.5pt]
  coordinates {(3,9) (1.2,1.44) (0.48,0.23) (0.19,0.04)};
\draw[-{Stealth[length=1.5mm]}, greenhl] (axis cs:3,9) -- (axis cs:1.2,1.44);
\draw[-{Stealth[length=1.5mm]}, greenhl] (axis cs:1.2,1.44) -- (axis cs:0.48,0.23);
\node[font=\tiny, text=greenhl!60!black, align=center] at (axis cs:-2.0,10.5) {converges\\quickly};

% --- too big
\nextgroupplot[title={\color{accent}$\alpha$ too large}]
\addplot[primary, line width=1pt, domain=-4.4:4.4, samples=80] {x^2};
\addplot[accent, only marks, mark=*, mark size=1.5pt]
  coordinates {(1.5,2.25) (-2.1,4.41) (2.94,8.64)};
\draw[-{Stealth[length=1.5mm]}, accent] (axis cs:1.5,2.25) -- (axis cs:-2.1,4.41);
\draw[-{Stealth[length=1.5mm]}, accent] (axis cs:-2.1,4.41) -- (axis cs:2.94,8.64);
\node[font=\tiny, text=accent!80!black, align=center] at (axis cs:0,11.5) {overshoots;\\diverges};
\end{groupplot}
\end{tikzpicture}}%
{The learning rate is the single most important hyperparameter you will tune. Too small
wastes compute; too large diverges. \chref{neuralnets} returns to this with real training
curves.}{learning-rate}

\begin{pitfall}
\begin{tight}
  \item \textbf{Computing distances on unscaled features.} The feature with the
        biggest numeric range silently becomes the only feature that matters.
  \item \textbf{Assuming a zero gradient means the best answer.} It means
        \emph{flat}. For the bowl-shaped losses of \chref{regression} that is the
        global minimum; for the neural networks of \chref{neuralnets} it need not
        be.
  \item \textbf{Treating the learning rate as unimportant.} It decides whether
        training works at all.
\end{tight}
\end{pitfall}

%---------------------------------------------------------------------
\begin{fourlenses}{Vectors, Distance and Optimisation}
\lensLA{Each student becomes a vector of engagement features. Euclidean distance
between two students answers ``who behaves like whom'', which is what drives peer
grouping and the clustering of \chref{unsupervised}. The scaling warning matters
acutely here: attendance on 0--100 will otherwise drown quiz counts on 0--10.}

\lensPM{Each completed task becomes a vector (estimated hours, actual hours,
number of dependencies, team size). Distance identifies the historical tasks most
similar to a new one, giving a defensible estimate --- ``the five most similar
past tasks took 11--16 hours'' --- rather than a guess.}

\lensIOT{A one-minute window of sensor readings becomes a vector. Distance from the
fleet's normal-operation vector is the simplest possible anomaly score, and it runs
in a few arithmetic operations --- which is why it can execute on the device
itself, at the edge (\chref{cloudedge}).}

\lensSEC{A network session becomes a vector (bytes in, bytes out, duration, packet
count, distinct ports). Cosine similarity to known attack signatures ignores
volume and compares \emph{shape}, so a slow, low-volume version of an attack still
resembles the fast one. Attackers routinely change magnitude; changing shape is
harder.}
\end{fourlenses}

%---------------------------------------------------------------------
\begin{lab}[Distance and Descent from First Principles]
\begin{lstlisting}[language=Python]
import numpy as np

# --- vectors and distance -------------------------------------------
a = np.array([92,  8, 12])   # engaged student
b = np.array([45,  3,  0])   # disengaged student

print("dot product :", a @ b)
print("norm of a   :", np.linalg.norm(a))
print("distance a-b:", np.linalg.norm(a - b))     # ~48.8

# Now scale each feature to 0-1 and recompute: the answer changes a lot,
# because attendance no longer dominates purely by having a bigger range.
lo, hi = np.array([0,0,0]), np.array([100, 10, 25])
a_s, b_s = (a - lo)/(hi - lo), (b - lo)/(hi - lo)
print("scaled distance:", np.linalg.norm(a_s - b_s))

# --- gradient descent on L(w) = w^2 ---------------------------------
w, alpha = 3.0, 0.3
for step in range(6):
    grad = 2 * w                 # dL/dw
    w = w - alpha * grad
    print(f"step {step+1}: w={w:7.4f}  L={w**2:8.5f}")

# Change alpha to 1.2 and run again -- watch the loss increase.
\end{lstlisting}
\end{lab}

%---------------------------------------------------------------------
\begin{checkpoint}
\begin{tightnum}
  \item Compute the Euclidean distance between $[3, 4]$ and $[0, 0]$.
  \item $L(w) = w^2$ and you are at $w = -2$. What is the gradient, and does
        gradient descent move $w$ left or right?
  \item Two documents have cosine similarity 0.95 but Euclidean distance 40. What
        does that combination tell you about them?
\end{tightnum}
\end{checkpoint}

%---------------------------------------------------------------------
\begin{chaptersummary}
\begin{tight}
  \item An observation is a vector; a dataset is an $n \times p$ matrix. Every
        algorithm consumes this shape.
  \item Euclidean distance measures dissimilarity; it is dominated by whichever
        feature has the largest scale, so scale first.
  \item Cosine similarity compares direction while ignoring magnitude --- the
        right choice for text.
  \item A derivative is a slope; a gradient is a vector of slopes pointing
        steeply uphill.
  \item Gradient descent is $w \leftarrow w - \alpha \nabla L$. The learning rate
        $\alpha$ decides between crawling, converging and diverging.
\end{tight}
\end{chaptersummary}

\begin{reviewq}
  \item \textbf{(MCQ)} The gradient of a loss function points in the direction of:
        (a)~steepest decrease (b)~steepest increase (c)~the minimum (d)~zero error.
  \item \textbf{(MCQ)} For a dataset with 500 rows and 12 columns, the matrix $X$
        has dimensions: (a)~$12 \times 500$ (b)~$500 \times 12$ (c)~$500 \times
        500$ (d)~$12 \times 12$.
  \item \textbf{(Short)} Compute $\|[6, 8]\|$ and the distance between $[6,8]$ and
        $[2,5]$.
  \item \textbf{(Short)} Explain in your own words why gradient descent
        automatically takes smaller steps as it approaches a minimum.
  \item \textbf{(Short)} Give one scenario in which cosine similarity is clearly
        the right choice and Euclidean distance is clearly wrong.
  \item \textbf{(Applied)} Perform two steps of gradient descent by hand on
        $L(w) = (w-4)^2$ from $w_0 = 0$ with $\alpha = 0.4$. (Hint: $\nabla L =
        2(w-4)$.) State the value of $L$ after each step.
  \item \textbf{(Applied)} A colleague computes distances between servers using
        features \{uptime in seconds, CPU load 0--1, error count 0--50\}. Explain
        precisely what will go wrong and how to fix it.
\end{reviewq}
